\documentclass[11pt]{article}
\usepackage[margin=1in]{geometry}
\usepackage{amsmath,amssymb}
\usepackage{booktabs}
\usepackage{array}
\usepackage{xcolor}
\usepackage[most]{tcolorbox}
\usepackage{enumitem}
\usepackage{multicol}
\usepackage[colorlinks=true,linkcolor=blue!50!black,urlcolor=blue!50!black,citecolor=blue!50!black]{hyperref}

\newtcolorbox{takeaway}[1][]{colback=yellow!12,colframe=orange!70!black,
  title={#1},fonttitle=\bfseries,boxrule=0.6pt,arc=2pt,left=4pt,right=4pt,top=3pt,bottom=3pt}
\newtcolorbox{keybox}[1][]{colback=blue!4,colframe=blue!40!black,
  title={#1},fonttitle=\bfseries,boxrule=0.6pt,arc=2pt,left=4pt,right=4pt,top=3pt,bottom=3pt}

\newcommand{\kl}{\ensuremath{\kappa_\lambda}}
\newcommand{\mhh}{\ensuremath{m_{HH}}}
\newcommand{\bbtt}{\ensuremath{b\bar b\tau\tau}}
\newcommand{\ptt}{\ensuremath{p_{\mathrm T}}}
\newcommand{\A}{\textsc{a}} % ATLAS tag
\newcommand{\Cms}{\textsc{c}} % CMS tag
\newcommand{\pA}[1]{\textcolor{gray}{\small[ATLAS p.\,#1]}}
\newcommand{\pC}[1]{\textcolor{gray}{\small[CMS p.\,#1]}}
\setlist{itemsep=2pt,topsep=3pt}

\newcommand{\priZ}{\textcolor{red!75!black}{\textbf{[P0]}}}
\newcommand{\priO}{\textcolor{orange!85!black}{\textbf{[P1]}}}
\newcommand{\priT}{\textcolor{gray!65!black}{\textbf{[P2]}}}
\newcommand{\eff}[1]{{\footnotesize\textcolor{gray!65!black}{(#1)}}}
\newcommand{\sref}[1]{ {\footnotesize\textcolor{gray!60!black}{[#1]}}}
\newcommand{\dw}{\textcolor{green!30!black}{\textbf{Done:}}~}
\newcommand{\stDone}{\textcolor{green!50!black}{\textbf{[\checkmark\,done]}}}
\newcommand{\stPart}{\textcolor{orange!90!black}{\textbf{[partial]}}}

\title{\vspace{-2.5em}New challenges for an NSBI measurement of \kl\ in HH$\to\bbtt$,\\
relative to the first NSBI measurement (ATLAS off-shell $H\to ZZ\to4\ell$)}
\author{\normalsize\itshape working note --- draft for discussion}
\date{20 June 2026}

\begin{document}
\maketitle
\vspace{-1.6em}

\begin{takeaway}[Scope \& one-line conclusion]
This compares the systematics and background handling of the \textbf{first full NSBI measurement}
(ATLAS off-shell $H\to ZZ\to4\ell$; note HIGG-2023-01 = arXiv:2412.01600, Sandesara \emph{et al.}) with what a
\textbf{\kl\ measurement in HH$\to\bbtt$} must do (baselined on CMS AN-18-121, Amendola \emph{et al.}). The NSBI
\emph{skeleton} transfers directly; the new difficulty is almost entirely in the \emph{inputs}: \bbtt\ has a far
messier final state, several \textbf{data-driven backgrounds with no per-event MC density} (QCD, jet$\to\tau$
fakes), an \textbf{irreducible same-final-state single-Higgs background}, an order of magnitude more (and larger)
\textbf{experimental systematics that deform the very observable (\mhh) carrying the \kl\ signal}, and a
\textbf{\kl-dependent theory uncertainty}. None of these are NSBI-specific --- they hit any method --- but several
sit awkwardly inside a per-event density-ratio framework in ways the clean, all-MC off-shell analysis never faced.
Page references to both notes are given inline.
\end{takeaway}

\section*{Studies to do}
\noindent{\small Concrete studies to de-risk and execute this measurement, grouped by theme. Priority:
\priZ\ blocking / do-first, \priO\ important, \priT\ nice-to-have; effort \eff{S $<$1\,wk, M weeks, L months}.
Trailing \sref{\S n} point to the relevant numbered section below. All abbreviations are defined in the
\textbf{Glossary} at the end. Status tags: \stDone\ $=$ completed and \stPart\ $=$ partially completed in the
preliminary phenomenological study (see the Progress box).}

\begin{keybox}[Progress --- preliminary phenomenological result (Ghosh, Klute \& Pan, 22 Jun 2026)]
{\small A first \bbtt\ NSBI study (code in \texttt{NSBI-pheno/dihiggs\_bbtautau}) has completed several items
below. \emph{It is explicitly phenomenological --- \textsc{powheg}$+$\textsc{delphes}, no experimental
systematics, fast-sim --- so the statistics and physics-method items are largely done, while the systematics,
in-situ-CR and data-driven-background items remain open.} \textbf{Headline:} an unbinned 10-feature NSBI likelihood
lifts the binned-\mhh\ secondary minimum near \kl$\,\sim\,$3--4 and \textbf{narrows the 68\% interval by
$\sim$an order of magnitude}, raising the \kl$=0$ rejection from $1.1\sigma$ (binned \mhh) to $\mathbf{2.75\sigma}$
at $450\,\mathrm{fb}^{-1}$ and to $\mathbf{7.1\sigma}$ at the HL-LHC ($3000\,\mathrm{fb}^{-1}$, vs $2.8\sigma$
binned); 500-toy pulls $\sim\mathcal N(0,1)$ for both fits (Figs.\,1--2, Table\,I). A focused test statistic (FTS)
further tightens the 95\% interval near the SM by 10--20\% at no extra training cost (Fig.\,3, Table\,II).
\textbf{Completed:} \stDone\ \#11, \#14, \#17. \textbf{Partial} (done without systematics, or by an alternative
route): \stPart\ \#1, \#4, \#7, \#13, \#15, \#16, \#18, \#20, \#24, \#25. Remaining items are systematics-,
CR-, and data-driven-background-related.}
\end{keybox}

\begin{keybox}[Critical path]
{\small Statistics-first: port the \textsc{jax}+\textsc{iminuit} fit (\#11) and fix $P_{\rm ref}$ (\#17), then build
the in-situ-CR likelihood (\#12) and the canonical ABCDnn flow with ExtendedABCD yields (\#1,\#2) in parallel; only
\emph{then} run Asimov closure (\#13) and Neyman coverage (\#14). The JES$\times\tau$ES/THU factorisation (\#5),
the coherent $\tau$ES$\to$SVfit chain (\#6), and the feature-ablation degeneracy study (\#18) must land
\emph{before} coverage so the morphing is trustworthy; honest GP morphing-uncertainty (\#9) is gated behind the
MC-stat ensemble (\#15). Multi-channel (\#23), the $k_{2V}$ tail (\#21), POI-correlations (\#22), and
pipeline/pretraining (\#25,\#26) proceed as P1/P2 follow-ups.}
\end{keybox}

\subsection*{Data-driven backgrounds}
\begin{enumerate}[leftmargin=2.6em]\small
\item \stPart\ \priZ\,\eff{L} \textbf{Canonical ABCDnn flow for QCD \& jet$\to\tau$ fakes} (yield/shape split, three-piece
uncertainty). MMD-trained NAF (continuous $(N_j,N_b)$ conditioning, identity warm-up) morphs MC $\to$ the
data-driven CR target; yield from ExtendedABCD/fake-factors, flow $\to$ per-event ratio $r_{\rm QCD}(x)$ via a
flow-vs-reference classifier; port the three four-top nuisances as per-event morphing; sweep target dim
2D$\to$5D$\to$10D. \dw VR closure $|1-\langle f\rangle|<0.05$, RMS $<0.15$; $r(x)\ge0$ integrating to the ABCD
yield; max stable dimension documented. \textcolor{orange!75!black}{\emph{\small(Prelim.: fake-$\tau$ and QCD are
already given a per-event density via parametrised mis-ID \emph{weights} on \textsc{powheg}/MG5 MC; the ABCDnn flow
$+$ closure on real data remains.)}} \sref{\S4}
\item \priZ\,\eff{M} \textbf{ExtendedABCD vs standard-ABCD yield closure for \bbtt\ QCD.} Standard ABCD failed
$\sim$18--19\% at high $N_j$ in the four-top analysis and the QCD-norm NP is already large; test the extended
estimator's closure across 6--8 orthogonal SRs binned in $(N_j,N_b)$. \dw bias $<7\%$, RMS $<12\%$ per
multiplicity bin; a safe-vs-extended region map. \sref{\S4}
\item \priO\,\eff{S} \textbf{Encode the one-sided QCD shape NP into smooth morphing.} CMS uses one alternative
template for both up/down; test mirrored / one-sided-exponential / half-Gaussian encodings so the poly-exp
morphing of $r_{\rm QCD}(x)$ does not oscillate. \dw matches the CMS shape-impact rank $\pm2$, no oscillation,
$<2\%$ \kl\ bias. \sref{\S4}
\item \stPart\ \priO\,\eff{L} \textbf{Quantify the binned-QCD/fakes vs fully-unbinned cost.} Asimov (\kl$=0,1,5$): full
unbinned vs MC-only NSBI with QCD/fakes as a fixed binned template; also bin the NSBI score at 5/10/20/50 to check
convergence. \dw width inflation $<15\%$, bias $<0.3\sigma$, monotonic convergence; a bin-vs-unbin fallback
decision tree. \sref{\S4,\,\S7}
\end{enumerate}

\subsection*{Systematics \& learned morphing}
\begin{enumerate}[resume,leftmargin=2.6em]\small
\item \priZ\,\eff{M} \textbf{Validate factorisation \& cross-terms: JES$\times\tau$ES and THU$_{HH}\times$\kl.} Both
JES and $\tau$ES feed \mhh\ (a cross-term absent in clean $4\ell$) and THU$_{HH}$ is POI-coupled. Test
$r_{\rm comb}/\prod_m r_m$ (median, KL) on held-out $\pm1\sigma$ MC; compare naive-lnN vs \kl-coupled THU in
Asimov. \dw factorisation residual $0.98$--$1.02$, KL $<0.02$; add a correction-NN if $>2\%$. \sref{\S5}
\item \priZ\,\eff{S} \textbf{Validate the coherent $\tau$ES$\to$MET$\to$SVfit$\to$\mhh\ recompute.} The
weight-builder must shift $\tau$ energy at object level and recompute MET$+$SVfit$+$\mhh\ end-to-end, not shift
\mhh\ directly. \dw matches the coherent truth to $<2\%$ per-event MAE; the shortcut bias quantified. \sref{\S3,\,\S5}
\item \stPart\ \priO\,\eff{M} \textbf{Validate poly-exp interpolation (\kl\ scan $+$ $\pm1\sigma$ anchors) under limited MC.}
Intermediate \kl\ (0.5, 1.5) vs reweighted truth; downsample $\pm1\sigma$ samples 100$\to$10\%. \dw MAE $<5\%$ down
to $\sim$50\% of MC, no gradient oscillation; minimum $\pm1\sigma$ sample size to request from generators. \sref{\S1,\,\S5}
\item \priO\,\eff{M} \textbf{Build \& profile a CARL/DCTR per-event parton-shower/generator nuisance} (the NSBI
upside). Train the per-event weight Pythia$\leftrightarrow$Herwig / ISR/FSR / VBF dipole-recoil; profile as one
Gaussian NP vs the binned envelope. \dw KS $p>0.1$, interval $\le$ envelope with correct coverage, generator
closure absorbed $<2\sigma$. \sref{\S5}
\item \priO\,\eff{M} \textbf{Honest GP/ensemble-spread morphing-uncertainty with coverage calibration.} Fit a GP to
per-event $r_{\rm true}-r_{\rm pred}$ residuals as a phase-space-dependent band (gated behind \#15). \dw 65--70\%
coverage, calibration slope $1\pm10\%$, Spearman $\rho>0.6$, narrower than a global band. \sref{\S5,\,\S7}
\item \priT\,\eff{M} \textbf{Absorb b-tag/$\tau$-ID SFs into the ratio, calibrated in situ} ($Z\to b\bar b$ /
VZ($c\bar c$)). Demote the POG SF program to a cross-check and shrink the NN count. \dw in-situ SF within 30\% of
the POG prior, \kl\ not degraded, $|\Delta\text{pull}(0\text{ vs }5)|<0.5\sigma$. \sref{\S5}
\end{enumerate}

\subsection*{Statistics, fit \& coverage}
\begin{enumerate}[resume,leftmargin=2.6em]\small
\item \stDone\ \priZ\,\eff{M} \textbf{Port the \textsc{jax}+\textsc{iminuit} profile fit with the
$\{$\kl$^2,$\kl$,1\}$ basis} (root dependency). Replicate ATLAS Sec.\,5: $P_{\rm ref}$, the
signal/interference/background mixture with \kl-dependent rate, $\pm1\sigma$ ratio classifiers $+$ poly-exp, profile
fit with Gaussian-constrained NPs; cross-check vs a binned Asimov. \dw converges, Hesse invertible, parabolic NLL;
unbinned-vs-binned width $<5\%$. \textcolor{green!45!black}{\emph{\small(Prelim.: unbinned extended likelihood
[Eq.\,3] with cached density ratios $+$ profile-NLL scan [Fig.\,1]; binned \mhh\ cross-check via \texttt{iminuit}.
The NP-profiling framework for the $\sim$100 experimental systematics is the remaining piece.)}} \sref{\S1,\,\S6}
\item \priZ\,\eff{M} \textbf{Represent the in-situ control regions (18 DY SFs, $t\bar t$-SF, QCD B/D) inside the
unbinned likelihood.} Add Poisson CR auxiliary terms sharing the same NPs; compare SR-only-free-NP vs SR$+$explicit
CRs. \dw reproduce the CMS in-situ SFs within 10--20\%, no double-counting, corr(SF,\kl)$<0.3$. \sref{\S4,\,\S6}
\item \stPart\ \priZ\,\eff{M} \textbf{Asimov \& closure: morphing bias under injected shifts and generator/Delphes
mismatch.} Inject known $\pm1\sigma$ on JES/$\tau$ES/QCD/MC-stat and re-fit; fit data from chain $X$ with a model
trained on $Y$ (Pythia$\leftrightarrow$Herwig; Delphes$\leftrightarrow$full-sim). \dw recover \kl\ to $<0.05$, pulls
$\sim N(0,1)$, sim-closure bias $<0.1$ (flag full-sim retraining if $>10\%$). \sref{\S5,\,\S6}
\item \stDone\ \priZ\,\eff{M} \textbf{Neyman construction \& coverage validation on MC-truth toys.} 1000$+$ toys; build the
belt at \kl$\in\{-1,0,1,3,5\}$. \dw empirical coverage $\ge67\%$, $|$bias$|<0.1$ at SM; any failure traced to a
specific NP/component. \textcolor{green!45!black}{\emph{\small(Prelim.: 500 toys at \kl$=1$ give pulls
$\sim\mathcal N(0,1)$ and calibrated 68\% intervals for both NSBI and binned [Fig.\,2]; extend to the full
\kl\ grid and re-run with systematics.)}} \sref{\S6}
\item \stPart\ \priO\,\eff{M} \textbf{MC-stat ensemble-spread nuisance: design, size, coverage, Barlow--Beeston
orthogonality.} Compare ensemble strategies (seeds / $k$-fold / bootstrap / \mhh-stratified) for the per-event
ratio. \dw C2ST $p>0.1$ on holdout, NP width $\approx1$ (BB-valid), \kl\ widens 5--15\% when unfrozen, corr $<0.3$
with physics NPs. \sref{\S5,\,\S6}
\item \stPart\ \priO\,\eff{S} \textbf{Interference-aware impact definition $+$ binned-NSBI fallback.}
Implicit-function-theorem impact at the minimum vs finite-difference and the CMS rank; histogram the per-event
log-ratio (5/10/20 bins) as an exact, learning-free cross-check. \dw matches the CMS rank $\pm2$; binned constraint
$\ge80\%$ of unbinned, monotonic convergence. \sref{\S6,\,\S7}
\end{enumerate}

\subsection*{Physics \& POI}
\begin{enumerate}[resume,leftmargin=2.6em]\small
\item \stDone\ \priZ\,\eff{S} \textbf{Define \& validate the reference hypothesis $P_{\rm ref}$.} Compare SM /
\kl-flat-mixture / pure-background / dedicated references; measure C2ST closure and how the
$\{$\kl$^2,$\kl$,1\}$ reconstruction degrades with $|$\kl$-$\kl$_{\rm ref}|$, including the thin high-\mhh\
$k_{2V}$ tail. \dw C2ST $p>0.05$ across the scan, $|\log r|<10$, $<5\%$ basis bias at extremes.
\textcolor{green!45!black}{\emph{\small(Prelim.: a single \emph{distant} reference \kl$=5$ was chosen; a
pooled-mixture reference gave 0.50 classifier accuracy [target--reference self-overlap] and was rejected;
$r_{\rm ref}\equiv1$; stable against reference-sample replacement.)}} \sref{\S1}
\item \stPart\ \priZ\,\eff{M} \textbf{Feature-ablation degeneracy-breaking: $\rho($\kl,JES/$\tau$ES$)$ and VIF vs feature
schedule.} \mhh$\to$10-feature$\to+$constituents$\to+$SVfit covariance$\to+$resolution handles; ablate to find the
critical features and the asymptotic $\rho$ floor. \dw reproduce $\rho\,1.000\to0.646$; a minimal set with
$\rho<0.7$, VIF$<1.5$; settle whether constituents/SVfit push $\rho<0.60$ or it plateaus.
\textcolor{orange!75!black}{\emph{\small(Prelim.: the \emph{degeneracy-break is demonstrated} --- the 10-feature
NSBI lifts the binned-\mhh\ secondary minimum near \kl$\,\sim\,$3--4 [Fig.\,1]; features selected by
total-variation distance $+$ integral-closure, stable under one-at-a-time low-TV removal. The
$\rho($\kl,JES/$\tau$ES$)$/VIF quantification needs the systematic variations.)}} \sref{\S2,\,\S5}
\item \priO\,\eff{S} \textbf{Source the \kl-differential THU$_{HH}$ theory band.} Compile the scale/PDF/$m_t$
uncertainty vs \kl$\in[-2,5]$ from the LHC HH WG / Powheg-NLO and parametrise it as a POI-coupled weight. \dw a
smooth band reproducing $+4/-18\%$ at the SM, ready for the factorisation study (\#5). \sref{\S2,\,\S5}
\item \stPart\ \priO\,\eff{M} \textbf{Validate the LO$\to$NLO \kl-reweighting that generates the $R(x|$\kl$)$ sources.}
Direct-NLO vs 2D-LO-reweight vs dedicated-NLO-reweight; propagate the difference through an Asimov fit. \dw
NLO-reweight reproduces direct to $<5\%$/bin and $<0.1\sigma$ \kl\ bias; the LO-only cost quantified.
\textcolor{orange!75!black}{\emph{\small(Prelim.: sidestepped --- signal generated \emph{directly} at NLO
[\textsc{powheg}] at \kl$\in\{0,1,5\}$ with a 3-point quadratic basis [the $\kappa_\lambda=2.45$ point was dropped
for morphing stability]; this study still applies if the CMS LO$+$reweighting samples are used instead.)}} \sref{\S2}
\item \priO\,\eff{M} \textbf{Measure $k_{2V}$ sensitivity in the high-\mhh\ VBF tail via unbinned NSBI.}
VBF-enriched ($M_{jj}>400$\,GeV), $k_{2V}$-morphed MC; Asimov-fit binned vs unbinned vs hybrid; cross-validate the
tail (train low-\mhh, test high). \dw NSBI $k_{2V}$ error 20--40\% smaller, tail-extrapolation bias $<10\%$, a
minimum viable bin granularity. \sref{\S2}
\item \priO\,\eff{L} \textbf{Map the rate-broken POI correlations ($\kappa_t\!\leftrightarrow$\kl,
$k_V\!\leftrightarrow k_{2V}$) and confirm the scope caveat.} Multi-POI grid; profile-likelihood correlation
contours; attribute information to shape vs rate. \dw a POI correlation matrix; NSBI lowers $\rho($\kl,$k_{2V})$ but
does \emph{not} improve the rate-broken pairs. \sref{\S2,\,\S7}
\item \priO\,\eff{M} \textbf{Build the multi-channel $+$ boosted simultaneous unbinned likelihood.} Per-channel vs
joint morphing for the three $\tau\tau$ channels $+$ boosted, with the CMS year/channel NP-correlation scheme. \dw
combined \kl\ width within 5\% of the best combination; no extra $>0.5\sigma$ pulls. \sref{\S1,\,\S3}
\end{enumerate}

\subsection*{ML, samples \& pipeline}
\begin{enumerate}[resume,leftmargin=2.6em]\small
\item \stPart\ \priO\,\eff{M} \textbf{Low-signal robustness: morphing calibration \& uncertainty honesty at
$\mathcal O(10$--$100)$ signal events.} Downsample the $\pm1\sigma$ sets and effective signal yield (50/100/250) and
retrain per level. \dw C2ST $p>0.05$ at the $\pm1\sigma$ floor, ensemble spread $<2\times$ the high-stat value;
flag the binned fallback if it fails. \sref{\S5,\,\S7}
\item \stPart\ \priT\,\eff{L} \textbf{Reproducible Perlmutter pipeline for the $\sim$2000-NN morphing ensemble.}
Containerise Sophon$+$\textsc{jax}$+$\textsc{iminuit}, a training manifest, resume-on-failure; benchmark per-NP
classifiers vs a single conditional morphing network. \dw 500 NNs in $<4$\,h ($\sim$1 GPU-day), $>99\%$ resume,
bit-identical reruns; the per-NP-vs-conditional trade-off documented. \sref{general}
\item \priT\,\eff{M} \textbf{Pretraining follow-up: signal-vs-$t\bar t$ foundation-model benchmark vs scratch}
(confirmatory). Scratch vs $t\bar t$-kinematics-pretrain vs HH-kinematics-pretrain; F-score \kl$=0$ vs 5 by \mhh\
region. \dw signal-vs-$t\bar t$ pretraining closes $\ge30\%$ of the FM deficit, else confirm scratch remains the
baseline. \sref{established null}
\end{enumerate}

\section{The two measurements and the shared NSBI skeleton}
\textbf{ATLAS off-shell} \pA{1,5}: measures the off-shell signal strength $\mu_{\text{off-shell}}$ (hence the Higgs
width) in $gg\to(H^*\to)ZZ\to4\ell$, using $139\,\mathrm{fb}^{-1}$. The physics is a \emph{coupling that morphs an
interference shape} (triangle--box in $ggZZ$). \textbf{CMS HH$\to\bbtt$} \pC{3,4}: searches for non-resonant HH via
ggF+VBF in \bbtt\ ($\mathcal B=7.3\%$), $137.6\,\mathrm{fb}^{-1}$, three $\tau\tau$ channels
($\tau_e\tau_h,\tau_\mu\tau_h,\tau_h\tau_h$, plus boosted), extracting \kl\ (and $k_{2V},k_V$). $\sigma_{ggF}^{HH}=
31.05^{+2.2}_{-5.0}\%$ fb, $\sigma_{VBF}^{HH}=1.726$ fb \pC{4}.

\noindent\textbf{What transfers (the skeleton).} The ATLAS recipe (Sec.\,5) is reusable wholesale: a
reference-hypothesis $P_{\mathrm{ref}}$ so all NNs estimate ratios to a fixed point; a search-oriented
\emph{mixture model} for signal/interference/background; nuisance parameters in both the rate and the per-event
density; a profile-likelihood test statistic in a custom \textsc{jax}+\textsc{iminuit} fit; and an
interference-aware impact definition \pA{89,90}. The off-shell interference structure
$(\mu-\sqrt\mu)\,S+\sqrt\mu\,\mathrm{SBI}+(1-\sqrt\mu)\,B$ \pA{81} maps onto the \kl\ morphing basis
$\{\kl^2,\kl,1\}$ (since $gg\to HH=\kl\cdot\mathrm{triangle}+\mathrm{box}$). \emph{The challenges below are about the
inputs to this skeleton, not the skeleton itself.}

\section{Physics / parameter-of-interest challenges}
\begin{itemize}
\item \textbf{The POI lives in a shape, and that shape is \mhh.} \kl\ sensitivity is concentrated in the \mhh\
  spectrum (triangle--box interference), strongest in the resolved regime $\mhh\in[250,400]$ GeV. CMS encodes the
  \kl\ dependence by \emph{reweighting} the signal MC --- a 2D $(\mhh,|\cos\theta^*|)$ LO reweighting and a
  dedicated NLO reweighting for the final extraction \pC{6}. In NSBI this becomes the $R(x|\kl)$ factor.
\item \textbf{Multiple correlated POIs / EFT directions.} Beyond \kl\ the analysis also probes $k_{2V},k_V$ and
  HEFT operators ($c_2,c_{2g},c_g$) \pC{4}; the off-shell measurement is effectively one POI. More POIs $\Rightarrow$
  a higher-dimensional morphing.
\item \textbf{$k_{2V}$ sensitivity lives in the high-\mhh\ VBF tail --- where binning hurts most.} The $HHVV$
  contact term grows with energy unless $k_{2V}=1$, so VBF-HH $k_{2V}$ sensitivity is concentrated in the
  high-\mhh\ tail; current analyses use only $\sim$4--5 bins there. This is precisely the regime where coarse
  binning discards the most information, so the unbinned NSBI gain is largest on the $k_{2V}$ (and the low-\mhh\
  \kl) shape directions. ($k_{2V}$ is now the best-measured ``inaccessible'' coupling --- the 2026 ATLAS+CMS
  combination gives $k_{2V}\in[0.73,1.3]$.)
\item \textbf{Interference impact ambiguity (shared, but worse).} Both have interference, so the
  finite-difference NP-impact definition is ambiguous; ATLAS replaces it with a local implicit-function-theorem
  impact \pA{90}. \kl\ needs the same fix.
\end{itemize}

\section{Final state and reconstruction}
The off-shell $4\ell$ final state is \emph{exquisitely clean}: four leptons, fully reconstructed, with tiny
energy-scale systematics and no jets/MET/$\tau$/$b$-tagging in the core. \bbtt\ is the opposite:
\begin{itemize}
\item \textbf{$b$-jets} (JES/JER, $b$-tagging) + \textbf{hadronic $\tau$} ($\tau$-ID, $\tau$ energy scale) +
  \textbf{MET} (the $\tau$ neutrinos). The $\tau\tau$ mass is \emph{not} fully reconstructed --- it requires an
  algorithm (SVfit); shifting $\tau$ energy shifts MET \emph{and} SVfit coherently \pC{93}.
\item \textbf{Three channels} with different background compositions and systematics, plus a boosted regime.
\end{itemize}
Consequence: the observable $x$ for \bbtt\ is softer and carries many more systematic ``handles'' than $4\ell$.

\section{Background handling --- the largest qualitative difference}
\begin{keybox}[The crux]
Off-shell backgrounds are \textbf{all MC-based and clean} (one dominant irreducible $q\bar q ZZ$ continuum, modelled
in simulation, plus EW). \bbtt\ backgrounds are a \textbf{zoo, several of them data-driven}, and \emph{data-driven
backgrounds have no per-event MC probability density} --- which is the object NSBI needs. This is the single
hardest conceptual mismatch between \bbtt\ and the off-shell NSBI template.
\end{keybox}

\noindent\textbf{ATLAS off-shell (Sec.\,2.5, Sec.\,5).} The dominant background is the irreducible $q\bar q\to
ZZ\to4\ell$ continuum, which \emph{interferes} with the signal and is carried in the mixture model as a component
$\mu_{q\bar qZZ}$ (three floating NPs) constrained \emph{in situ} by a dedicated $n_{\text{jets}}$ discriminant in a
control region \pA{81}. Plus a small EW component. Everything is MC, so every process has a per-event density the
NSBI ratios can target.

\noindent\textbf{CMS \bbtt\ (Sec.\,7).} A heterogeneous set, with the two largest estimated \emph{from data}:
\begin{itemize}
\item \textbf{QCD multijet --- data-driven ABCD} \pC{61,62}. Four regions in (charge $\times$ $\tau$-isolation);
  shape from region C (data minus other-MC), yield from $C\times B/D$. The QCD \emph{normalisation} uncertainty
  ranges \textbf{5--100\%} by category (Table 30, \pC{92}); its \emph{shape} uncertainty (alternative B-vs-C
  templates, symmetrised) is \textbf{one of the highest-ranking nuisances} in the most sensitive
  $\tau_h\tau_h$ channel \pC{91,96,97}.
\item \textbf{Drell--Yan ($Z$+jets) --- MC normalised by 18 data-driven scale factors} \pC{63}. A simultaneous fit
  of $M(\mu\mu)$ across \textbf{18 control regions} (3 $b$-jet multiplicities $\times$ 6 $Z$-\ptt\ bins), giving 18
  normalisation SFs ($\sim$0.6--1.7) \emph{and} 18 shape templates (scaled $\pm$) propagated with their covariance
  \pC{65,94}.
\item \textbf{$t\bar t$ --- MC shape + data-driven normalisation SF} \pC{79}. A dedicated CR (inverted elliptical
  $(\mhh)$ mass cut, Eq.\,22) with $<7\%$ contamination; \textsc{combine} fit with Barlow--Beeston gives
  $\text{ttSF}\approx0.91$--$0.99$ \pC{81,82}; the shape is taken from MC (no shape dependence observed).
\item \textbf{Jet$\to\tau$ fakes --- data-driven} \pC{94,95}. $>93\%$ of the events with a jet faking a $\tau$;
  fake-rate uncertainty per region (barrel 24/13/11\%, endcap 28/18/25\% by year, Table 33), applied as an
  event weight.
\item \textbf{Single Higgs --- irreducible, same final state} \pC{82}. ggH, VBF$H$, $VH$, $t\bar tH$ with $H\to
  b\bar b$ or $\tau\tau$ mimic the signal; modelled from MC. (Note: $t\bar tH$/$VH$ couplings make this a
  \emph{coupling-dependent} background.)
\item \textbf{Di-/tri-boson, $W$+jets, single-$t$, $V$+2jets --- MC} \pC{82}.
\end{itemize}

\begin{takeaway}[Why this is a new, NSBI-specific difficulty]
\begin{enumerate}[leftmargin=1.4em]
\item \textbf{No density for data-driven backgrounds.} QCD (ABCD) and jet$\to\tau$ fakes are estimated as
  \emph{yields/weights from data control regions}, not as simulated events with a per-event probability. NSBI's
  per-event likelihood ratio $P_X(x)/P_{\mathrm{ref}}(x)$ is undefined for them. Folding an ABCD/fake estimate into
  an \emph{unbinned} per-event likelihood is an open problem the off-shell analysis never had to solve (its one big
  background was MC). Options (all imperfect): treat QCD/fakes only in a binned auxiliary term; build a surrogate
  density from the control-region data; or restrict NSBI to the MC-modelled part and bin the rest.
\item \textbf{The QCD shape systematic is a top-ranked nuisance} \pC{91} --- and it is intrinsically
  \emph{asymmetric/ill-defined} (one alternative template used for both up and down), forcing a symmetrisation
  hack \pC{96,97}. Carrying that into a per-event morphing is non-trivial.
\item \textbf{Single Higgs is an irreducible, coupling-dependent background} in the same final state --- it must be
  modelled \emph{and} morphed consistently with the signal couplings, with its own (correlated) theory systematics.
\item \textbf{Many background-normalisation NPs are themselves data-fit results} (18 DY SFs with a covariance,
  ttSF, QCD $B/D$) --- these are constrained \emph{in situ} in the binned CMS fit; reproducing that constraint
  inside an unbinned NSBI fit requires the CRs to be represented in the per-event likelihood too.
\end{enumerate}
\end{takeaway}

\subsection*{A workaround: flow-based data-driven backgrounds (ABCDnn)}
The ``no per-event density'' problem is the sharpest objection to NSBI here, but it is \emph{not} a dead end ---
there is a demonstrated solution. \textbf{ABCDnn} (Choi, Lim, Oh, arXiv:2008.03636) generalises the ABCD method
with a \emph{neural autoregressive flow} that learns the full multi-dimensional CR$\to$SR transformation,
conditioned on the discriminating observables. Where standard ABCD gives only a scalar SR yield
($N_D=N_B N_C/N_A$), ABCDnn gives a full \emph{differential} background distribution $p_{\rm QCD}(x)$ --- which can
then be turned into a density ratio $r_{\rm QCD}(x)=p_{\rm QCD}(x)/p_{\rm ref}(x)$ by training a classifier between
flow-generated events and the pooled reference, so it enters the NSBI likelihood \emph{on the same footing} as the
MC-derived ratios.

\noindent\textbf{The concrete production example} (S.~An, Dutta/Paulini, CMU PhD thesis 2025, Ch.\,5): the CMS
\emph{all-hadronic four-top} search, Run-2 ($137\,\mathrm{fb}^{-1}$, 13 TeV) + Run-3 ($62\,\mathrm{fb}^{-1}$,
13.6 TeV), reaching an \textbf{observed (expected) significance of $3.95\,(1.65)\,\sigma$}, $\mu=2.7^{+1.1}_{-0.9}$.
Its single most important design choice is a \textbf{yield/shape split} of the dominant $t\bar t$+jets and QCD
backgrounds:
\begin{itemize}
\item \emph{Yield} $\leftarrow$ \textbf{ExtendedABCD}: standard ABCD in $(N_{\rm jets},N_{b})$ \emph{fails} here
  (the factorisation breaks at high multiplicity --- off by $\sim$18--19.5\% in closure); adding two extra
  sidebands restores it to $\sim$1--7\% (their $\hat F_D=(F_BF_C/F_A)^2(F_X/F_BF_Y)$, Table 5.1--5.2). The lesson:
  \emph{get the normalisation from a robust closed-form estimator with enough sidebands, not from the flow.}
\item \emph{Shape} $\leftarrow$ \textbf{ABCDnn} (a neural autoregressive flow): it learns a transformation that
  \emph{morphs the $t\bar t$ MC into the data-driven $t\bar t$+QCD estimate} (target $=$ data $-$ signal $-$ minor
  bkg per CR), trained by minimising \textbf{Maximum Mean Discrepancy} and \emph{conditioned on the CR position
  $(N_{\rm jets},N_b)$}, then applied to $t\bar t$ MC in the SR. Notably it never uses QCD MC (too few SR events) ---
  the flow turns pure-$t\bar t$ MC into an effective $t\bar t$+QCD mixture.
\end{itemize}
Three training tricks they found necessary (worth copying): \textbf{loss-weighted CR sampling} (sample harder CRs
more often, $\propto e^{\rm loss}$, so every region is learned uniformly); \textbf{continuous $(N_{\rm jets},N_b)$
conditioning} instead of one-hot labels, so the flow \emph{interpolates smoothly} across CRs --- this is what makes
the CR$\to$SR \emph{extrapolation} well-behaved; and a \textbf{warm-up} (pre-train near identity, RealNVP-style) for
stable convergence.

\noindent\textbf{Application to \bbtt:} the \textbf{jet$\to\tau$ fakes} are the natural target --- a fake-factor
method is conceptually an ABCD with $\tau$-ID inversion as the second axis, so an ABCDnn conditioned on
$(p_{\rm T}^\tau,\eta^\tau)$ could learn the full fake-$\tau$ kinematic shape instead of the current coarse 2D
fake-factor binning; the QCD multijet is the other target; and for $t\bar t$ the flow can serve as a data-driven
shape cross-check. Following the four-top lesson, keep the \textbf{yield/shape split}: take the fake-$\tau$/QCD
\emph{normalisation} from the existing fake-factor/ABCD machinery (which CMS already trusts, \pC{61,94}) and use
ABCDnn only for the differential \emph{shape}, then convert that shape to a per-event ratio $r(x)$ for the
likelihood. A fully realised \bbtt-NSBI would use \emph{ABCDnn shapes (with ABCD/fake-factor yields) for QCD/fakes,
MC ratios for the MC-modelled components ($Z$+jets, $t\bar t$, single-$H$, diboson), and one unbinned NSBI
likelihood combining both} --- removing the main caveat that NSBI ``depends on MC where MC is wrong.''

\noindent\textbf{But the generic caveats remain:} (i) \emph{extrapolation distance} --- the CR$\to$SR map strains for
tight SRs far from the sidebands in tagger/isolation space (multiple overlapping CRs validate it but cost stats);
(ii) \emph{uncertainty propagation} --- the four-top analysis does \emph{not} bootstrap the flow; it splits the
uncertainty into three data-driven nuisances (thesis Sec.\,6.1, Table 6.1): the \emph{finite-CR-statistics} piece
is a $\ln$N normalisation NP propagated through the ABCD formula (``statistics of extendedABCD,'' 4--26\%); the
\emph{normalisation-modelling} piece is a $\ln$N NP from a validation-region (VR) data-vs-prediction closure
(weighted-mean deviation $1-\langle f\rangle$ $+$ weighted RMS of $f=N_{\rm data}/N_{\rm pred}$, 6--44\%); and the
\emph{ABCDnn shape} uncertainty is a shape NP implemented as a \emph{linear shift of the discriminant}, with the
shift size chosen by \emph{minimising a quantitative data/pred-agreement metric} in the VR (floor 1\%, correlated
across $H_T$, uncorrelated across years). The whole uncertainty-\emph{estimation} procedure is itself validated in a
dedicated orthogonal ``closure-test region'' that checks pred$+$syst covers truth in the SR. \emph{Caveat for NSBI:}
this VR-closure shape nuisance (``shift the discriminant'') is a \emph{binned-fit} construct; in an unbinned NSBI it
becomes yet another per-event morphing nuisance on the ratio $r(x)$ --- i.e.\ the same learned-morphing problem,
now applied to the data-driven component; (iii) \emph{signal contamination in CRs}
--- if signal leaks into the control region the flow learns ``background $+$ signal'' and biases the SR low,
requiring iterative CR definitions with signal subtraction (this caveat is \emph{mild} for \bbtt: the HH signal
is tiny in the fake-$\tau$/QCD control regions, just as the SM four-top signal was negligible in its CRs); (iv) \emph{no theory handle} --- the flow gives a shape
but no ME$+$PS$+$PDF decomposition, so generator-level theory uncertainties on the background cannot be propagated
and are replaced by a data-driven CR$\to$SR closure systematic.

\noindent\textbf{The one caveat specific to using ABCDnn \emph{for NSBI} (the real gap).} In the four-top thesis
the flow models only a \emph{low-dimensional} target --- the event-level BDT discriminant plus $H_T$ (two
variables) --- because the downstream fit is a \emph{binned template} fit of that discriminant. NSBI needs the
opposite: the \emph{full per-event density} over the whole feature vector, so the data-driven shape can become a
per-event ratio $r(x)$. So the four-top result proves flows can deliver a data-driven background \emph{shape of a
summary observable}; extending that to the full-dimensional per-event density an unbinned NSBI consumes is a genuine
step beyond what is demonstrated --- a higher-dimensional flow, harder to train and validate. The conservative
intermediate is to let ABCDnn produce the data-driven shape of the \emph{NSBI discriminant} (the per-event
log-ratio summed over components), i.e.\ treat QCD/fakes as a binned component of an otherwise-unbinned likelihood
--- exactly the ``keep QCD binned, NSBI-fy the MC components'' compromise.

\noindent\emph{Net: ABCDnn downgrades the data-driven-background problem from ``open'' to ``solved-with-caveats''
for a \emph{binned} discriminant shape; the fully-unbinned version is a worthwhile but unproven sub-project for a
\bbtt-NSBI pilot.}

\section{Systematics handling}
\textbf{What is genuinely shared} (equal for both, and equally fallible): the \emph{sources and $\pm1\sigma$ sizes}
of the NPs; the \emph{object-level propagation} (shift the object and all dependents --- e.g.\ $\tau\Rightarrow$
MET$\Rightarrow$SVfit --- and recompute the discriminant, \pC{93}); the \emph{factorisation over NPs} and
$\pm1\sigma$ \emph{interpolation}; \emph{Gaussian-constrained profiling}; and the \emph{\kl-dependent theory}
treatment. ATLAS builds the per-event systematic weight as $P_X(x,\alpha)=w_X(x,\alpha)\,P_X(x,0)$ with
$w_X(x,\alpha)=\prod_m w_X(x,\alpha_m)$, each single-NP weight interpolated poly-exp between $\pm1\sigma$ classifier
estimates \pA{83,84}. CMS does the per-\emph{bin} analogue (re-run the DNN on varied inputs $\to$ template).

\medskip\noindent\textbf{What is harder for \kl/\bbtt:}
\begin{enumerate}[leftmargin=1.4em]
\item \textbf{\kl--systematic shape degeneracy.} The dominant experimental systematics (JES, $\tau$ES, MET, signal
  parton shower) deform \emph{the same \mhh\ shape} that carries \kl. In $4\ell$ the big systematics (theory) are
  largely \emph{separable} from the signal shape; in \bbtt\ they are not. Decorrelation is therefore impossible
  (it would erase the signal), forcing the ``model-it'' route --- so the accuracy of the morphing
  $\Delta(x|\alpha)$ \emph{directly sets} the \kl\ uncertainty. (A Fisher-information toy gives, schematically,
  a $\sim$13$\times$ systematic inflation from $\mhh$ alone collapsing to $\sim$1.3$\times$ with the full feature
  set --- the degeneracy is broken by dimensionality, not avoided.)
\item \textbf{An order of magnitude more, and larger, experimental NPs} \pC{90--97}: JES split into a reduced set
  of \textbf{11 sources} with a year-correlation scheme \pC{95}; \textbf{$\tau$ES} \ptt-dependent with 4 decay-mode
  components \pC{94}; \textbf{$b$-tag} 5 components (hf, lf, hfstats, lfstats, cferr) \pC{96}; \textbf{DeepTau}
  vsJet (5 \ptt\ bins) + vsEle (2 $\eta$ bins) \pC{95,96}; trigger SFs, JER, PU-jet-ID, L1 prefiring, $e/\mu\to\tau$
  fake ES. The off-shell analysis has comparatively few, small (lepton) experimental NPs. Each NP is another
  per-event weight to learn ($\sim$2000 NNs already in off-shell $=2\times$NP$\times$process; \bbtt\ grows this).
\item \textbf{Data-driven background systematics that don't fit the density-ratio picture}: the QCD norm (5--100\%)
  and high-ranking QCD shape \pC{91,92,96,97}, the 18 DY SF shape/norm templates \pC{94}, the jet$\to\tau$ fake-rate
  uncertainties (11--28\%, Table 33) \pC{95}, the VBF dipole-recoil PS uncertainty (20--70\%, Table 32) \pC{93}.
  These are template/weight constructs, not deformations of an MC density.
\item \textbf{\kl-dependent theory uncertainty (POI$\times$systematic entanglement).} The HH scale/PDF/$m_t$
  uncertainties are folded into a normalisation NP \texttt{THU\_HH} $=+4\%/-18\%$ at the SM point that becomes a
  \emph{\kl-dependent} uncertainty across the scan \pC{91}. The off-shell $\mu$ does not entangle with its theory
  uncertainty this way. Capturing it requires non-factorisable (POI-coupled) morphing.
\item \textbf{Factorisation is more likely to break.} $w_X=\prod_m w_X(\alpha_m)$ assumes NPs are independent
  sources \pA{83}; but JES and $\tau$ES \emph{both} feed \mhh\ ($=m(b\bar b+\tau\tau)$), so a genuine cross-term
  exists in \bbtt\ that is benign in clean $4\ell$.
\item \textbf{MC-statistical uncertainty of the learned ratios.} ATLAS propagates it as one NP via the
  NN-ensemble spread \pA{86,87}; CMS uses Barlow--Beeston per bin. With far smaller \emph{systematic-variation}
  samples in \bbtt, this term is more important (the nominal MC, $\approx$500k--750k/\kl-point, is ample for
  training; the binding constraint is the least-populated $\pm1\sigma$ sample).
\end{enumerate}

\medskip\noindent\textbf{Two systematics that NSBI actually handles \emph{better} than templates.} Worth stating
the upside too: (a) \emph{parton-shower / generator} uncertainties (Pythia vs.\ Herwig, ISR/FSR) are a natural fit
--- a \textbf{CARL/DCTR}-style reweighting network learns the per-event weight between alternative showers, which is
then profiled as a nuisance; this is strictly more powerful than the current ``shape envelope over generator
variations'' and is the in-likelihood cousin of the RS3L robust-representation idea. (b) $b$-tag / $\tau$-ID scale
factors can be \emph{absorbed} into the density ratio and calibrated in situ against a standard candle (the
$VZ(c\bar c)$/$Z\to b\bar b$ peak trick), turning the POG SF program into an external cross-check rather than the
primary constraint. So the systematics ledger is not all cost: \emph{the morphing-heavy ones (JES/$\tau$ES on
\mhh) are the burden; the modelling ones (shower, flavour-tagging) are where NSBI's per-event treatment is an
asset.}

\section{Statistics and inference}
\begin{itemize}
\item \textbf{Data signal yield is tiny and intrinsic.} $\sigma\times\mathcal B\sim2.5$ fb $\Rightarrow$
  $\mathcal O(10$--$100)$ selected signal events (Run-2 to HL-LHC), under large backgrounds --- vs the off-shell
  tail at $\sim$15\% of $\sigma$. This is unchanged by MC and keeps the measurement systematics-comparable. (It is
  also the regime where NSBI's per-event information helps \emph{most}.)
\item \textbf{MC/training stats are ample} ($\approx$500k--750k per \kl\ point), so the morphing/ratio NNs are
  well-trained --- the ``under-determined morphing'' worry applies to the \emph{systematic-variation} samples, not
  the nominal.
\item \textbf{Validation \& coverage.} CMS inherits decades-validated \textsc{combine} (asymptotics + toys +
  Barlow--Beeston). NSBI needs a custom \textsc{jax}/\textsc{iminuit} fit with an explicit Neyman construction and
  ratio-quality diagnostics (C2ST-like) built and validated from scratch \pA{64,89}.
\end{itemize}

\section{Synthesis: could the NSBI \kl\ result be \emph{worse} than the CMS fit?}
Reading what CMS \emph{actually} does removes a strawman: CMS does \textbf{not} bin \mhh; it bins a
\textbf{multivariate DNN score} (shape systematics propagated by re-running the DNN on varied inputs, \pC{93}). So
CMS already carries the same multi-dimensional systematic information NSBI does --- \textbf{NSBI is not ``exposed to
more systematics'' than the real CMS baseline.} Against CMS-as-done, the comparison reduces to:
\begin{itemize}
\item \textbf{What NSBI buys:} \emph{unbinning} (no information lost to binning the score) and a \emph{\kl-optimal}
  per-event ratio (CMS's DNN is one fixed discriminant, not optimal at every \kl). These are the off-shell
  $\sim$2$\times$ / FAIR-Universe $\sim$20\% gains.
\item \textbf{What NSBI costs:} a \emph{learned} morphing layer (NSBI must \emph{train} the per-event systematic
  weight; CMS \emph{counts} its re-run-DNN template, exact up to MC stat) and the \emph{custom coverage
  validation}. The danger of ``worse than CMS'' is the \emph{learned layer being mis-calibrated / under-validated}
  --- producing a tighter-but-mis-covering interval --- not a heavier systematics burden.
\item \textbf{The genuinely new \bbtt\ problem is the data-driven backgrounds} (Sec.\,4): if QCD/fakes cannot be
  given a faithful per-event density, the unbinned NSBI may have to bin them anyway, partially eroding its
  advantage in exactly the most sensitive channel.
\end{itemize}
Protections (so the worst case is ``$\approx$ CMS,'' not ``worse and unknown''): honest morphing uncertainty
(GP-style); explicit coverage validation; and \textbf{binning the NSBI discriminant as a built-in cross-check}
(the binned fit is a coarsening of NSBI, so a regression is detectable and you fall back).

\section{Summary of new challenges}
\begin{table}[h]\centering\small
\begin{tabular}{@{}p{0.26\textwidth}p{0.32\textwidth}p{0.34\textwidth}@{}}
\toprule
 & \textbf{ATLAS off-shell $4\ell$} & \textbf{\kl\ in HH$\to\bbtt$ (new challenge)} \\
\midrule
POI & $\mu_{\text{off-shell}}$ (1 POI) & \kl\ ($+k_{2V},k_V$, EFT); shape in \mhh \\
Final state & 4 leptons, fully reco, clean & $b$-jets+$\tau_h$+MET; SVfit; 3 channels \\
Backgrounds & one MC $q\bar qZZ$ (interfering), CR-constrained & QCD (ABCD), DY (18-CR SFs), $t\bar t$ (CR-SF), \textbf{jet$\to\tau$ fakes}, single-$H$, diboson \\
Data-driven bkg in NSBI & none needed & \textbf{QCD \& fakes have no per-event density} \\
Experimental systematics & few, small (leptons) & many, large (JES$\times$11, $\tau$ES, $b$-tag$\times$5, DeepTau, \ldots) \\
Syst.\ vs signal shape & largely separable & \textbf{JES/$\tau$ES deform \mhh\ (degenerate with \kl)} \\
Theory--POI coupling & none & \textbf{\kl-dependent \texttt{THU\_HH} $+4/-18\%$} \\
Factorisation of NPs & safe & JES$\times\tau$ES cross-term in \mhh \\
Signal yield & $\sim$15\% of $\sigma$ & $\mathcal O(10$--$100)$ events (fb-scale) \\
MC training stats & --- & ample ($\approx$500--750k/\kl); watch $\pm1\sigma$ samples \\
Stats framework & custom (built) & must port the custom fit; CMS has \textsc{combine} \\
\bottomrule
\end{tabular}
\end{table}

\begin{takeaway}[Bottom line]
The NSBI \emph{method} ports cleanly from off-shell to \kl. The new challenges are in the \emph{inputs}, and rank
roughly: (1) \textbf{data-driven backgrounds} (QCD, jet$\to\tau$ fakes) with no per-event density --- the most
awkward fit with NSBI, but \emph{solved-with-caveats} by flow-based estimation (\textbf{ABCDnn}, Sec.\,4); (2) the
\textbf{\kl--systematic shape degeneracy} (JES/$\tau$ES on \mhh) plus a much larger experimental-systematics
surface, which makes \emph{morphing accuracy + honesty} the limiting factor; (3) the \textbf{\kl-dependent theory}
and \textbf{factorisation cross-terms} (non-factorisable morphing); (4) \textbf{irreducible single-Higgs}. None are
NSBI-specific \emph{except} the data-driven-background mismatch and the learned-morphing layer; the rest are equally
CMS's problems, handled with the same NPs. Two scope caveats: NSBI's gain is in the \emph{\kl/$k_{2V}$ parameter
constraint} (shape directions), \emph{not} the SM-HH observation significance (rate-saturated, luminosity-set);
and the genuine per-event edge is on \emph{shape/correlated-operator} directions, not the rate-broken
$\kappa_t\!\leftrightarrow\!\kl$ / $k_V\!\leftrightarrow\!k_{2V}$ ones.
\end{takeaway}

\section*{Glossary of abbreviations \& jargon}
\noindent\small Every acronym and piece of jargon used in this note, in plain language. Symbols/couplings first,
then abbreviations alphabetically.\normalsize

\subsection*{Symbols \& couplings}
{\footnotesize\setlength{\parindent}{0pt}\setlength{\parskip}{2.5pt}
\begin{multicols}{2}
$\kappa_\lambda$ (\kl) --- Higgs trilinear self-coupling strength relative to the SM ($=1$ in the SM); the
parameter of interest. Enters gg$\to$HH through the \emph{triangle} diagram and interferes with the \emph{box}.

$k_{2V}$ --- strength of the quartic $VVHH$ contact coupling ($=1$ in SM); accessed in the high-\mhh\ VBF tail.

$k_V$ --- single Higgs--vector-boson coupling strength ($VVH$, $=1$ in SM).

$\kappa_t$ --- top-quark Yukawa coupling strength ($=1$ in SM); enters HH via the box and triangle.

\mhh\ --- invariant mass of the di-Higgs system; the observable carrying the \kl\ interference shape.

$\cos\theta^*$ --- production angle of the Higgs in the HH rest frame; the second variable in CMS's 2D \kl\ reweighting.

$H_T$ --- scalar sum of jet transverse momenta.

$M_{jj}$ --- invariant mass of the two VBF ``tagging'' jets.

$N_j,\,N_b$ --- number of jets / number of $b$-tagged jets (the ABCD(nn) conditioning variables).

$p_{\mathrm T}$ --- transverse momentum (momentum in the plane transverse to the beam).

$P_{\rm ref}$ --- reference hypothesis: the single fixed model all NSBI classifiers estimate density ratios \emph{relative to}.

$r(x)$ --- per-event likelihood (density) ratio --- the quantity NSBI learns and the fit consumes.

$\rho$ --- (Pearson/Spearman) correlation coefficient; here between \kl\ and a systematic.

$\sigma$ --- standard deviation / Gaussian width (``$\pm1\sigma$'' $=$ a one-standard-deviation systematic variation); also denotes a cross-section.

triangle / box --- the two gg$\to$HH amplitudes: the \kl-sensitive \emph{triangle} (via an $s$-channel Higgs) and the \kl-independent \emph{box}; their interference shapes \mhh.

GeV, TeV --- giga-/tera-electronvolt (energy/mass units).
\end{multicols}}

\subsection*{Abbreviations \& jargon}
{\footnotesize\setlength{\parindent}{0pt}\setlength{\parskip}{2.5pt}
\begin{multicols}{2}
\textbf{ABCD} --- data-driven background estimate from four regions defined by two (assumed-independent) variables; the signal region $D=(B\times C)/A$.

\textbf{ABCDnn} --- a normalizing-flow generalization of ABCD that morphs simulated background into the data-driven background \emph{shape}, conditioned on extra variables.

\textbf{AN} --- (CMS) Analysis Note: internal documentation (e.g.\ AN-18-121).

\textbf{anchors} --- the fixed reference points (a \kl\ grid, or the $\pm1\sigma$ systematic variations) between which a smooth morphing is interpolated.

\textbf{Asimov dataset} --- a single artificial ``data $=$ expectation'' dataset used to compute \emph{expected} sensitivity without generating toys.

\textbf{ATLAS, CMS} --- the two general-purpose LHC experiments.

\textbf{AUC} --- area under the ROC curve; a classifier separation metric ($0.5=$ no separation).

\textbf{Barlow--Beeston} --- standard treatment of finite-MC-statistics uncertainty via per-bin nuisance parameters.

\textbf{BDT} --- boosted decision tree.

\textbf{binned / unbinned} --- a fit that histograms an observable into template bins vs.\ one that uses each event's per-event likelihood directly (no binning). The CMS baseline is binned; NSBI is unbinned.

\textbf{bootstrap} --- resampling with replacement to estimate a statistical uncertainty (e.g.\ of a trained network).

\textbf{CARL / DCTR} --- the classifier-reweighting family: train a classifier to learn a per-event reweighting between two models. CARL $=$ Calibrated Approximate Ratio of Likelihoods; DCTR $=$ Deep Classifier-based ReweighTing.

\textbf{closure (test)} --- a check that a method reproduces a known truth (e.g.\ predicted $=$ true background in a validation region).

\textbf{C2ST} --- Classifier Two-Sample Test: train a classifier to distinguish two samples; AUC $\approx0.5$ means they are indistinguishable (good closure).

\textbf{Combine} --- CMS's standard binned maximum-likelihood fitting framework (HistFactory-style; asymptotic formulae plus toys, with Barlow--Beeston MC-stat handling).

\textbf{CR / SR / VR} --- Control / Signal / Validation Region.

\textbf{decay mode} --- the specific hadronic $\tau$ decay channel (e.g.\ 1-prong, 3-prong, with/without $\pi^0$); $\tau$ES has per-decay-mode components.

\textbf{DeepTau} --- CMS deep-learning hadronic-$\tau$ identification discriminant.

\textbf{Delphes} --- a fast, parametrised detector simulation used for quick studies before full simulation.

\textbf{DNN} --- deep neural network.

\textbf{DY} --- Drell--Yan ($Z/\gamma^*\to\ell\ell$) background.

\textbf{EFT / SMEFT} --- (Standard-Model) Effective Field Theory: parametrize new physics as higher-dimensional operators.

\textbf{ES} --- energy scale (as in JES, $\tau$ES).

\textbf{EW} --- electroweak.

\textbf{ensemble (spread)} --- training several networks (different seeds/folds); their spread estimates the MC-statistical uncertainty, propagated as one nuisance.

\textbf{ExtendedABCD} --- ABCD augmented with extra sidebands to correct residual correlation between the two defining variables (a more accurate yield).

\textbf{extrapolation (CR$\to$SR)} --- applying a transfer/calibration learned in a control region to the signal region; its validity is the key data-driven-background caveat.

\textbf{fake-factor} --- data-driven estimate of jet$\to\tau_h$ misidentification: a transfer factor measured in a $\tau$-anti-isolated CR and applied to the SR.

\textbf{Fisher information} --- quantifies how sharply an observable constrains a parameter; here used to show that added dimensions shrink the \kl--systematic degeneracy.

\textbf{flavour-tagging} --- identifying a jet's originating quark ($b$/$c$/light) with an ML discriminant ($b$-tagging is the $b$-jet case).

\textbf{F-score} --- harmonic mean of precision and recall; here a separation proxy.

\textbf{ggF / VBF} --- gluon--gluon fusion / vector-boson fusion (the two HH production modes).

\textbf{ggH, VBF\,H, VH, $t\bar tH$ (single Higgs)} --- the single-Higgs production modes (gluon fusion / vector-boson fusion / associated with a vector boson / with a top pair); the irreducible same-final-state background (distinct from ggF/VBF for HH).

\textbf{GP} --- Gaussian Process: a non-parametric model giving predictions with calibrated uncertainty bands.

\textbf{HH} --- di-Higgs (Higgs-pair) production; \textbf{HHVV} $=$ the quartic $VVHH$ contact vertex (set by $k_{2V}$).

\textbf{Hesse / Hessian} --- matrix of second derivatives of the NLL; its inverse gives parameter covariances (\textsc{minuit}'s \textsc{hesse}).

\textbf{HL-LHC} --- High-Luminosity LHC.

\textbf{ID} --- identification (e.g.\ $\tau$-ID, $b$-ID).

\textbf{iminuit} --- Python interface to the \textsc{minuit} minimizer used for the profile fit.

\textbf{interference} --- the quantum cross-term between two amplitudes (here triangle$\times$box in HH); it carries the \kl\ shape information.

\textbf{ISR / FSR} --- initial-/final-state radiation (a parton-shower systematic).

\textbf{JAX} --- an autodiff / accelerated-array library used to build the differentiable likelihood.

\textbf{JER / JES} --- jet energy resolution / jet energy scale (systematics).

\textbf{KL divergence} --- Kullback--Leibler divergence: a measure of difference between two probability distributions.

\textbf{KS test} --- Kolmogorov--Smirnov goodness-of-fit test (a $p$-value for ``same distribution'').

\textbf{LHC} --- Large Hadron Collider.

\textbf{lnN} --- log-normal: the usual constraint shape for a \emph{normalization} nuisance parameter.

\textbf{LO / NLO} --- leading / next-to-leading order in perturbative QCD.

\textbf{MAE} --- mean absolute error.

\textbf{MC} --- Monte Carlo (simulated events).

\textbf{ME} --- matrix element.

\textbf{MET} --- missing transverse energy (momentum imbalance, here mainly from neutrinos).

\textbf{ML} --- machine learning.

\textbf{mixture model} --- writing the total per-event density as a weighted sum of signal/interference/background components (the NSBI fit model).

\textbf{MMD} --- Maximum Mean Discrepancy: a distribution-matching loss (used to train ABCDnn).

\textbf{morphing} --- smoothly deforming a per-event density or weight as a function of a parameter (\kl\ or a nuisance), built by interpolating between anchors.

\textbf{NAF} --- Neural Autoregressive Flow: a flexible normalizing-flow density model (ABCDnn's network).

\textbf{Neyman construction} --- the frequentist procedure for building confidence intervals with correct coverage.

\textbf{NLL} --- negative log-likelihood.

\textbf{NN} --- neural network.

\textbf{NP} --- nuisance parameter: a non-POI (systematic) parameter that is profiled in the fit.

\textbf{NSBI / SBI} --- (Neural) Simulation-Based Inference: inference from learned per-event likelihood ratios instead of binned templates.

\textbf{PAS} --- (CMS) Physics Analysis Summary, a short public document.

\textbf{PDF} --- parton distribution function (separately, also ``probability density function'').

\textbf{Perlmutter} --- the NERSC GPU supercomputer used for training.

\textbf{pivoting / decorrelation} --- training a classifier to be \emph{insensitive} to a nuisance (statistically independent of it).

\textbf{POG} --- (CMS) Physics Object Group; provides standard object calibrations / scale factors.

\textbf{POI} --- parameter of interest (here \kl).

\textbf{poly-exp} --- polynomial$\times$exponential interpolation used between systematic ($\pm1\sigma$) anchors.

\textbf{Powheg} --- an NLO matrix-element event generator.

\textbf{profile (likelihood)} --- maximizing the likelihood over the nuisance parameters at each value of the POI.

\textbf{PS / PU} --- parton shower / pile-up (additional simultaneous $pp$ collisions).

\textbf{Pythia, Herwig} --- two parton-shower/hadronization generators; their difference gauges shower-modeling systematics.

\textbf{QCD} --- quantum chromodynamics; ``QCD multijet'' $=$ the data-driven multijet background.

\textbf{RealNVP} --- a specific (identity-initializable) normalizing-flow architecture.

\textbf{resolution} --- the smearing/precision of a measured quantity (cf.\ JER for jets); ``resolution handles'' $=$ input features encoding it.

\textbf{RMS} --- root-mean-square.

\textbf{RS3L} --- Re-Simulation-based Self-Supervised Learning (a representation-robustness method).

\textbf{Run-2 / Run-3} --- LHC data-taking periods (Run-2: 2016--18; Run-3: 2022--25).

\textbf{SF} --- scale factor: a data/MC correction applied to reconstructed objects ($b$-tag, $\tau$-ID, \ldots).

\textbf{sideband} --- a region adjacent to (but outside) the signal region, used to estimate backgrounds; the extra regions in ExtendedABCD are sidebands.

\textbf{SM} --- Standard Model.

\textbf{Sophon} --- the (Sophon-AK4) jet ``foundation model'' backbone fine-tuned in the ML-pipeline studies.

\textbf{SVfit} --- algorithm that reconstructs the $\tau\tau$ invariant mass from the visible decay products $+$ MET.

\textbf{THU$_{HH}$} --- theory uncertainty on the HH cross-section/shape; \kl-dependent ($+4/-18\%$ at the SM point).

\textbf{$\tau$ES (tauES)} --- $\tau$ energy scale (a systematic).

\textbf{toy / toy MC} --- synthetic pseudo-datasets drawn from the model, used for coverage and sensitivity tests.

\textbf{trigger} --- the online selection deciding which collision events to record; carries its own efficiency scale factor.

\textbf{VH / VZ} --- associated production of a Higgs / $Z$ with a vector boson.

\textbf{VIF} --- variance inflation factor: how much a parameter's variance grows due to correlation with others (a degeneracy metric).

\textbf{WG} --- working group (e.g.\ the LHC HH WG).

\textbf{ZZ, 4$\ell$} --- a $Z$-boson pair / four charged leptons (the ATLAS off-shell final state).
\end{multicols}}

\begin{thebibliography}{9}\small
\bibitem{atlas} ATLAS Collaboration (J.~Sandesara \emph{et al.}), \emph{Measurement of the Off-Shell Higgs Boson in
  $H\to ZZ\to4\ell$ using Simulation-Based Inference}, ATLAS note HIGG-2023-01 (2024) = arXiv:2412.01600,
  Rep.\ Prog.\ Phys.\ 88 (2025). Page refs herein are to the note. Key: mixture model \& systematics Sec.\,5
  (pp.\,80--90); per-event factorised weights (pp.\,83--84); MC-stat NP (pp.\,86--87); JAX fit \& impact
  (pp.\,89--90).
\bibitem{cms} CMS Collaboration (C.~Amendola \emph{et al.}), \emph{Search for Higgs boson pairs produced through
  gluon and vector boson fusion in the \bbtt\ final state with Run-II data}, CMS AN-18-121 (2022). Key: signal
  reweighting Sec.\,3.1 (p.\,6); background estimation Sec.\,7 (pp.\,61--85; QCD ABCD pp.\,61--62, DY 18 CRs
  pp.\,63--78, $t\bar t$ CR-SF pp.\,79--82); systematics Sec.\,9 (pp.\,90--97).
\bibitem{abcdnn} S.~Choi, J.~Lim, H.~Oh, \emph{Data-driven Estimation of Background Distribution through Neural
  Autoregressive Flows (ABCDnn)}, arXiv:2008.03636 (2020). Production application: CMS all-hadronic four-top search,
  Run-2+3 (ExtendedABCD yield + ABCDnn shape for $t\bar t$+jets/QCD; obs.\ $3.95\,\sigma$, $\mu=2.7^{+1.1}_{-0.9}$)
  --- S.~An (advisors M.~Paulini, V.~Dutta), Carnegie Mellon PhD thesis (2025), CMU KiltHub \#29950337;
  MMD-trained NAF, conditioned on $(N_{\rm jets},N_b)$, modelling the BDT$+H_T$ shape. No public CMS-PAS located.
\bibitem{bahl} A.~Bahl, T.~Plehn, N.~Schmal, \emph{global SMEFT fit of four di-boson processes with per-event
  likelihoods}, arXiv:2509.05409 (2025) --- the $\sim$2$\times$ per-event gain on correlated-operator directions.
\bibitem{comb} ATLAS+CMS HH combination, arXiv:2602.23991 / CMS-PAS-HIG-25-014 (early 2026):
  $k_{2V}\in[0.73,1.3]$, $\kl\in[-0.71,6.1]$, $\mu_{HH}<2.5$ (95\% CL).
\bibitem{gkp} A.~Ghosh, M.~Klute, J.~Pan, \emph{Focused Neural Simulation-Based Inference for the Higgs
  Self-Coupling in the $b\bar b\tau\tau$ Channel}, working note (22 Jun 2026); code:
  \texttt{NSBI-pheno/dihiggs\_bbtautau}. Phenomenological \textsc{powheg}$+$\textsc{delphes} study; signal
  $gg\to HH$ at NLO with $\kappa_\lambda\in\{0,1,5\}$ (three-point quadratic morphing basis); reference
  $\kappa_\lambda=5$; ten event-level observables; full background model with weighted fake-$\tau$; unbinned
  extended likelihood vs binned \mhh, plus a focused test statistic. Results: \kl$=0$ rejection $2.75\sigma$
  ($450\,\mathrm{fb}^{-1}$) $\to 7.1\sigma$ (HL-LHC); 68\% interval $\sim$10$\times$ tighter than the binned fit;
  500-toy coverage validated. \emph{No experimental systematics (phenomenological).}
\end{thebibliography}

\end{document}
